\section{Introduction}
\label{sec:introduction}

Computer numerical control (CNC) machining underpins precision manufacturing in aerospace, automotive, medical, and defense sectors, where dimensional deviations as small as 0.05~mm can render safety-critical parts non-functional~\cite{nist2023,graves2023cyberphysical}.
As manufacturing systems adopt Industry~4.0 architectures, cloud-based CAM, and remote monitoring, they inherit the cybersecurity vulnerabilities of general-purpose IT while retaining the physical consequences unique to operational technology (OT).
NIST SP 800-82 Rev.~3 identifies CNC controllers as critical OT assets requiring integrity verification~\cite{nist2023}, and the Cybersecurity Manufacturing Innovation Institute (CyManII) has prioritized CNC attack detection as a national research need~\cite{cymanii2022}.

The canonical cyber-physical attack modifies the G-code program (the motion and machine-control commands defining the manufacturing process) either in transit or at rest on the controller's file system~\cite{yampolskiy2018security,belikovetsky2017dr0wned}.
Even simple modifications (replacing a coordinate value, swapping a motion command, altering a feed rate) can produce parts that pass visual inspection but fail under load~\cite{belikovetsky2017dr0wned,zeltmann2016manufacturing}, as documented by the MITRE ATT\&CK for ICS framework under the ``Impair Process Control'' tactic~\cite{mitre2021attack}.

Existing CNC monitoring performs \emph{process classification} rather than \emph{integrity verification}.
Classification approaches assign sensor signals to known toolpath categories but cannot verify that the physically executed process matches the commanded program.
An adversary who modifies G-code to produce a slightly altered but physically valid toolpath (e.g., shifting a coordinate by 0.1~in or substituting one operation for another) generates sensor signals that remain within the distribution of normal machining, evading classification-based detectors entirely.

We propose a \emph{reconstruction-mismatch} approach to CNC anomaly detection.
A model trained to reconstruct G-code token sequences from multi-modal sensor signals learns the conditional distribution $p(\mathbf{y} \mid \mathbf{z})$ mapping sensor-derived latent representations $\mathbf{z}$ to G-code tokens $\mathbf{y}$.
When an attacker modifies the G-code, the sensor signals still reflect the \emph{original} physical process.
The decoder therefore reconstructs the correct tokens with high probability while assigning low probability to the substituted tokens.
The discrepancy between reconstruction and claimed G-code provides the detection signal.

This approach extends reconstruction-based anomaly detection~\cite{an2015variational,chalapathy2019deep} in three domain-specific ways.
First, the reconstruction target is a \emph{structured token sequence} (G-code), enabling language-model scoring methods (NLL, rank-based metrics, perplexity).
Second, multi-modal sensor input (accelerometers, gyroscopes, magnetometers, audio, environmental sensors, and electrical measurements across six sensor boards) constrains the reconstruction space far more than any single modality.
Third, the G-code grammar partitions attacks into \emph{grammar-violating} (trivially detectable by parsing) and \emph{grammar-compliant} (the hard, security-relevant case) categories, focusing evaluation on attacks that require learned representations.

The deployable \emph{unsupervised} detector reaches AUROC 0.803 on command-swap attacks (V7 NLL), 0.859 on small coordinate perturbations (V7 rank scoring), and 0.949 on cross-operation swaps (V8 NLL).
Sensor-only baselines, in contrast, lack a genuine detection signal against these G-code integrity attacks.
A supervised XGBoost classifier raises small-coordinate detection to 0.962 when labeled attacks are available; we report this as an upper bound rather than as the deployable result.
A per-token confidence analysis, an adaptive-adversary evaluation, and a calibration study characterize the limits of the approach.

The contributions of this paper are:

\begin{enumerate}
    \item \textbf{Multi-scale reconstruction-mismatch framework.} We combine a frozen multi-modal encoder (MM-DTAE-LSTM, 110 sensor channels) with two autoregressive decoder scales, short-context (V7) and multi-line (V8), and with rank-based scoring and ensemble methods, to detect grammar-compliant G-code attacks from sensor signals alone.

    \item \textbf{Supervised upper bound.} A supervised XGBoost classifier trained on features combining encoder memory, NLL statistics, and rank scores achieves AUROC 0.962 on small coordinate attacks when labeled attacks are available; we present this as an upper bound, not the deployable system. An unsupervised grid-searched ensemble only matches the best individual scorer per attack type (in-sample AUROC up to 0.859).

    \item \textbf{Baseline inadequacy.} We show that sensor-only baselines have no genuine detection signal for G-code integrity attacks. Because the attack leaves the physical process and every sensor channel unchanged, a \emph{paired} sensor-space comparison is exactly at chance, and the near-chance unpaired baselines we report deviate from 0.50 only through window-selection artifacts, not detection. The detection signal must instead come from comparing sensor observations against the \emph{claimed} program, which is what reconstruction-mismatch does.

    \item \textbf{Characterization of detection limits.} A per-token confidence analysis shows that detectability rises smoothly with the decoder's confidence; confident tokens (54\% of positions) reach per-token AUROC 0.89 while low-confidence tokens are near chance. We evaluate a first-order adaptive adversary that drives single-score detection to near chance (AUROC 0.54). We further document calibration failures (temperature scaling degrades ECE for the multi-line decoder) and identify a position confound affecting coordinate-attack evaluation.
\end{enumerate}

Code and trained models are available at [URL redacted for review].
